\chapter{Games of Incomplete Information}\label{ch:b2:bayesian-games}

In a sealed-bid auction each bidder knows what the object is worth to her but not what it is worth to the
others. Game theory predicts how much below her value she should bid in a first-price auction, where the
winner pays her own bid: with values drawn evenly between zero and a maximum, a bidder facing two rivals should
bid two thirds of her value, and one facing five rivals five sixths. The chart compares that prediction with
the bids of eighteen subjects in a laboratory experiment.

\begin{openingchart}{Fitted bid per unit of value of each of eighteen subjects in laboratory first-price
  auctions with values uniform between 0 and 30 dollars, bidding against 2 rivals and against 5. The dashed
  lines are the risk-neutral equilibrium predictions, 2/3 and 5/6.}{Bajari and Horta\c{c}su (2003), Table~1,
  on the experiments of Dyer, Kagel and Levin (1989)}
\begin{tikzpicture}
\begin{axis}[omrecord, xmin=0.3, xmax=18.7, ymin=0.6, ymax=1.0, xlabel={Subject},
    ylabel={Bid per dollar of value}, xtick={1,3,...,17}, width=0.85\linewidth, height=0.45\linewidth,
    legend style={at={(0.5,-0.25)}, anchor=north, legend columns=2,
    /tikz/every even column/.append style={column sep=10pt}}]
  \addplot[only marks, mark=*, mark size=2pt, omDef] table[x=subject, y=slope3]
    {figdata/bachelor-2/14-bayesian-games/out/bids.csv};
  \addplot[only marks, mark=square*, mark size=1.8pt, omThm] table[x=subject, y=slope6]
    {figdata/bachelor-2/14-bayesian-games/out/bids.csv};
  \addplot[omDef, dashed, domain=0.3:18.7, samples=2, forget plot] {2/3};
  \addplot[omThm, dashed, domain=0.3:18.7, samples=2, forget plot] {5/6};
  \legend{3 bidders, 6 bidders}
\end{axis}
\end{tikzpicture}
\end{openingchart}

% ledger: b2:bh2003-design, b2:bh2003-t1
% figdata: bids_summary (above3 18, above6 14, at6 1, below6 3, mean3 0.811, mean6 0.867)
With two rivals, every one of the eighteen subjects bid more than the prediction: on average 81 cents per dollar
of value against 67. With five rivals the bids rose, as predicted, and came much closer to it: fourteen subjects
bid above 5/6, one at it and three below, with an average of 87 cents.

So far every player has known the others' payoffs. Bidders do not; nor does a firm know its rival's costs, or an
employer a job applicant's ability. This chapter extends \omterm{def:b2:static-games:nash-equilibrium}{Nash equilibrium} to games in which players hold private
information, applies it to auctions, and shows why a bidder who does not reason about the information revealed
by winning tends to regret it.

\begin{recall}
From \cref{ch:b2:static-games,ch:b2:dynamic-games}: \omterm{def:b2:static-games:nash-equilibrium}{Nash equilibrium}, dominance and \omterm{def:b2:dynamic-games:spe}{subgame-perfect equilibrium}.
From \cref{ch:b2:expected-utility}: expected utility and \omterm{def:b2:expected-utility:risk-aversion}{risk aversion}. From the Year 1 volume: Cournot
competition.
\end{recall}

\begin{notation}[Symbols of this chapter]\label{not:b2:bayesian-games:symbols}
$\theta_i$ is player $i$'s \omterm{def:b2:bayesian-games:type}{type} and $p(\cdot)$ a probability. In auctions, $n$ is the number of bidders, $v_i$ bidder $i$'s value, $b_i$
her bid and $r$ a \omterm{def:b2:expected-utility:rra}{coefficient of relative risk aversion}. In the common-value model, $V$ is the common value, $s_i = V + \varepsilon_i$
bidder $i$'s signal and $\bar\varepsilon$ the largest possible error.
\end{notation}

\section{Types and beliefs}\label{sec:b2:bayesian-games:types}

\begin{definition}[Type]\label{def:b2:bayesian-games:type}
A player's \emph{type}\index{type (of a player)} is the private information she holds when she chooses: her
value for an object, her cost, her ability. Her payoffs may depend on her own type and on the others' types.
\end{definition}

\begin{definition}[Common prior]\label{def:b2:bayesian-games:common-prior}
Players share a \emph{common prior}\index{common prior} when the profile of \omterm{def:b2:bayesian-games:type}{types} is drawn from one probability
distribution that every player knows. Each player then forms her beliefs about the others' \omterm{def:b2:bayesian-games:type}{types} by conditioning
that distribution on her own \omterm{def:b2:bayesian-games:type}{type}.
\end{definition}

\begin{definition}[Bayesian game]\label{def:b2:bayesian-games:bayesian-game}
A \emph{Bayesian game}\index{Bayesian game} consists of players, their actions, their possible \omterm{def:b2:bayesian-games:type}{types}, a \omterm{def:b2:bayesian-games:common-prior}{common
prior} over profiles of \omterm{def:b2:bayesian-games:type}{types}, and payoffs that depend on the actions and the \omterm{def:b2:bayesian-games:type}{types}. A strategy is a function that
assigns an action to each of a player's \omterm{def:b2:bayesian-games:type}{types}.
\end{definition}

Conditioning uses Bayes' rule: the probability of $A$ given $B$ is $p(A \mid B) = p(B \mid A)\,p(A)/p(B)$. If 30 per cent of
fishing grounds are rich, and a survey reports ``rich'' for 80 per cent of rich grounds and 20 per cent of poor ones,
a report of ``rich'' raises the probability that a ground is rich to $0.24/(0.24 + 0.14) = 0.63$
(\cref{exo:b2:bayesian-games:7}). When \omterm{def:b2:bayesian-games:type}{types} are independent, as in most of this chapter, a player's own \omterm{def:b2:bayesian-games:type}{type} tells
her nothing about the others', and her beliefs are the prior itself.

\section{Bayesian Nash equilibrium}\label{sec:b2:bayesian-games:bne}

\begin{definition}[Bayesian Nash equilibrium]\label{def:b2:bayesian-games:bne}
A \emph{Bayesian Nash equilibrium}\index{Bayesian Nash equilibrium} is a profile of strategies in which each
\omterm{def:b2:bayesian-games:type}{type} of each player chooses an action that maximises her expected payoff, given her beliefs about the others'
\omterm{def:b2:bayesian-games:type}{types} and the others' strategies.
\end{definition}

The key step is that a player must predict what every \omterm{def:b2:bayesian-games:type}{type} of her rival would do, because she does not know which
\omterm{def:b2:bayesian-games:type}{type} she faces. Take two firms competing in quantities with demand $p = 12 - Q$. Firm 1's unit cost is 2 and
known to both. Firm 2's cost is 4 or 2, with probability one half each, and only firm 2 knows it. A high-cost
firm 2 best-responds with $q_2^H = (12 - 4 - q_1)/2$, a low-cost one with $q_2^L = (12 - 2 - q_1)/2$. Firm 1 best-responds to the
\emph{expected} output of firm 2, $q_1 = (12 - 2 - \tfrac12 q_2^H - \tfrac12 q_2^L)/2$. Solving the three equations,
\[
q_1 = 3.667, \qquad q_2^H = 2.167, \qquad q_2^L = 3.167.
\]
With full information, a high-cost firm 2 would face a firm 1 producing 4 and would produce 2; here it produces more,
because firm 1, unsure, produces less than it would against a known high-cost rival. A low-cost firm 2 produces less
than its full-information output 3.333, for the opposite reason.

\section{Auctions with private values}\label{sec:b2:bayesian-games:auctions}

\begin{definition}[Independent private values]\label{def:b2:bayesian-games:private-values}
In an auction with \emph{independent private values}\index{independent private values} each bidder's value for
the object is her \omterm{def:b2:bayesian-games:type}{type}, known only to her, drawn independently from a distribution that all bidders know, and
unaffected by what the others know.
\end{definition}

\begin{definition}[First-price sealed-bid auction]\label{def:b2:bayesian-games:first-price}
In a \emph{first-price sealed-bid auction}\index{first-price sealed-bid auction} every bidder submits one bid
without seeing the others'; the highest bidder wins and pays her own bid.
\end{definition}

\begin{definition}[Second-price sealed-bid auction]\label{def:b2:bayesian-games:second-price}
In a \emph{second-price sealed-bid auction}\index{second-price sealed-bid auction} every bidder submits one
bid without seeing the others'; the highest bidder wins and pays the second-highest bid.
\end{definition}

\begin{proposition}[Bidding one's value]\label{prop:b2:bayesian-games:truthful}
In a \omterm{def:b2:bayesian-games:second-price}{second-price sealed-bid auction} with private values, bidding one's value is a weakly dominant strategy.
\end{proposition}

\begin{proof}
Fix the others' bids and let $m$ be the highest of them. A bidder with value $v$ who bids $v$ wins when $v > m$, earning $v
- m > 0$, and loses otherwise, earning 0. Bidding $b > v$ changes the outcome only when $v < m < b$: she then wins and earns
$v - m < 0$. Bidding $b < v$ changes it only when $b < m < v$: she then loses a gain of $v - m > 0$. Her bid fixes whether she
wins, never what she pays, so no other bid ever does better than $v$, and each does worse for some bids of the others.
\end{proof}

In a first-price auction bidding one's value guarantees a payoff of zero, so bidders shade their bids. With values
uniform on $[0, 1]$, suppose the others bid a fixed share $k$ of their values. A bidder with value $v$ who bids $b$ beats one
rival with probability $p(k v_j < b) = b/k$, and all $n - 1$ with probability $(b/k)^{n-1}$. She maximises $(v - b)(b/k)^{n-1}$. The
first-order condition, $(n - 1)(v - b) = b$, gives $b = (n - 1)v/n$, which is of the form assumed, with $k = (n - 1)/n$.

\begin{proposition}[First-price bids with uniform values]\label{prop:b2:bayesian-games:first-price-bids}
In a first-price auction with $n$ risk-neutral bidders whose values are independent and uniform on $[0, 1]$, each bidding
$b(v) = (n - 1)v/n$ is a \omterm{def:b2:bayesian-games:bne}{Bayesian Nash equilibrium}.
\end{proposition}

\begin{proof}
Given the others' strategies, the derivation shows that $(n - 1)v/n$ maximises the bidder's expected payoff among bids
in $[0, k]$: the objective is zero at $b = 0$ and at $b = v$, and it rises then falls between them. A bid above $k$, the highest bid any
rival makes, wins for sure, but less than $k$ also wins for sure and costs less. So no \omterm{def:b2:bayesian-games:type}{type} gains by deviating.
\end{proof}

\begin{omfigure}
\begin{tikzpicture}[econ]
\begin{axis}[omrecord, grid=none, xmin=0, xmax=1.05, ymin=0, ymax=1.05, width=0.5\linewidth,
    height=0.5\linewidth, xlabel={Value $v$}, ylabel={Bid}, xtick={0,0.5,1}, ytick={0,0.5,1},
    legend style={at={(1.05,0.5)}, anchor=west}]
  \addplot[black, thick, domain=0:1, samples=2] {x};
  \addplot[omDef, thick, domain=0:1, samples=2] {x/2};
  \addplot[omThm, thick, dashed, domain=0:1, samples=2] {2*x/3};
  \addplot[omMeth, thick, dashdotted, domain=0:1, samples=2] {5*x/6};
  \legend{Second price (any $n$), First price $n = 2$, First price $n = 3$, First price $n = 6$}
\end{axis}
\end{tikzpicture}

{\small Equilibrium bids with values uniform on $[0, 1]$. In the second-price auction each bidder bids her value. In the
first-price auction she bids $(n - 1)/n$ of it: she shades her bid less when there are more rivals, since the highest of
them is likely to be close to her.}
\end{omfigure}

A bidder who is averse to risk shades less: a higher bid lowers her gain if she wins but makes winning more likely,
and she values that certainty. With \omterm{def:b2:expected-utility:crra}{constant relative risk aversion} $r$ the equilibrium bid becomes $(n - 1)v/(n - r)$
(Bajari and Horta\c{c}su 2003, following Cox, Smith and Walker 1988), so a measured slope of bids on values reveals $r$.
That is how the opening chart's slopes were estimated.

\section{Revenue equivalence}\label{sec:b2:bayesian-games:revenue}

Which format should a seller choose? In the first-price auction the winner pays her own, shaded, bid; in the
second-price she pays a rival's unshaded one. The effects cancel.

\begin{definition}[Revenue equivalence theorem]\label{def:b2:bayesian-games:revenue-equivalence}
The \emph{revenue equivalence theorem}\index{revenue equivalence theorem} is the result that, with risk-neutral
bidders whose values are independent draws from one distribution, every auction format in which the bidder with the
highest value wins and a bidder with the lowest possible value expects to pay nothing gives the seller the same
expected revenue.
\end{definition}

\begin{theorem}[Revenue equivalence]\label{thm:b2:bayesian-games:revenue}
Under the assumptions of the definition, the first-price and \omterm{def:b2:bayesian-games:second-price}{second-price sealed-bid auctions} give the same expected
revenue, and so does every other format satisfying them.
\end{theorem}

\admitted

The proof is given in the Year 3 volume. For uniform values the two revenues can be computed directly. The highest of
$n$ values uniform on $[0, 1]$ has expectation $n/(n + 1)$ and the second-highest $(n - 1)/(n + 1)$. The first-price auction yields the
winner's bid, $(n - 1)/n$ times the highest value, whose expectation is $(n - 1)/(n + 1)$; the second-price yields the
second-highest value, also $(n - 1)/(n + 1)$ (\cref{exo:b2:bayesian-games:11}). The listing checks this by simulation, and adds
bidders who are averse to risk.

\omcode{code/bachelor-2/14-bayesian-games/listings/auctions.py}{1}{19}{Revenue of first-price and second-price
auctions, simulated with values uniform between 0 and 1.}\label{lst:b2:bayesian-games:auctions}

\omoutput{code/bachelor-2/14-bayesian-games/outputs/auctions.txt}{Equal expected revenues; the first-price
revenue varies less, and rises when bidders are averse to risk.}\label{out:b2:bayesian-games:auctions}

% code: listings/auctions.py
The two formats raise the same revenue on average but not the same risk for the seller: the first-price revenue
varies about half as much, since it depends on one value and a shading that falls as the value rises. \omterm{def:b2:expected-utility:risk-aversion}{Risk aversion}
breaks the equivalence: it raises bids in the first-price auction and leaves bidding one's value optimal in the
second-price, so the first-price auction raises more.

\begin{model}{Auction with independent private values}
\assumes Risk-neutral bidders, values drawn independently from a known distribution, equilibrium play.
\predicts Bidding one's value in the second-price auction; bids of $(n - 1)v/n$ in the first-price auction with uniform
values; equal expected revenue.
\tested Laboratory auctions with assigned values (Dyer, Kagel and Levin 1989, as analysed by Bajari and Horta\c{c}su
2003).
\breaks Bids above the prediction in first-price auctions, most with few bidders, which \omterm{def:b2:expected-utility:risk-aversion}{risk aversion} can explain;
with common values, winning bids above the object's value (the \omterm{def:b2:bayesian-games:winners-curse}{winner's curse}).
\end{model}

\begin{record}[Bidding above the prediction]
% ledger: b2:bh2003-design, b2:bh2003-t1
In the experiments of Dyer, Kagel and Levin (1989), three groups of six business students bid for an object whose
value to each of them was drawn between 0 and 30 dollars. Each submitted two bids, one for an auction with 3 bidders and
one with 6, and a coin decided which applied. Bids were above the risk-neutral prediction: by 2.55 dollars on average,
in absolute value, with 3 bidders and by 0.94 with 6. Estimates of the \omterm{def:b2:expected-utility:risk-aversion}{risk aversion} that would explain the bids with 3
bidders are near 0.5, as in other auction experiments, but the bids with 6 imply much less \omterm{def:b2:expected-utility:risk-aversion}{risk aversion} in the same
people, which \omterm{def:b2:expected-utility:risk-aversion}{risk aversion} alone cannot account for (Bajari and Horta\c{c}su 2003).
\end{record}

\section{Common values and the winner's curse}\label{sec:b2:bayesian-games:common}

\begin{definition}[Common-value auction]\label{def:b2:bayesian-games:common-value}
In a \emph{common-value auction}\index{common-value auction} the object is worth the same amount $V$ to every bidder,
but no bidder knows $V$; each holds a private signal correlated with it, such as an estimate of the oil under a tract.
\end{definition}

Suppose bidder $i$'s signal is $s_i = V + \varepsilon_i$, with the errors $\varepsilon_i$ independent and uniform on $[-\bar\varepsilon, \bar\varepsilon]$. On its own a
signal is unbiased: the expected value given $s_i$ is $s_i$. But a bidder who wins with the highest bid holds, if all bid
alike, the highest signal, the one with the largest error. The highest of $n$ such errors has expectation
$\bar\varepsilon\,(n - 1)/(n + 1)$ (\cref{exo:b2:bayesian-games:12}), so
\[
E[V \mid s_i = s \text{ and } s_i \text{ is the highest}] = s - \bar\varepsilon\,\frac{n - 1}{n + 1}.
\]

\begin{definition}[Winner's curse]\label{def:b2:bayesian-games:winners-curse}
The \emph{winner's curse}\index{winner's curse} is the tendency of the winner of a \omterm{def:b2:bayesian-games:common-value}{common-value auction} to have
overestimated the object's value, so that she pays more than it is worth, when bidders fail to allow for the
information that winning reveals.
\end{definition}

\begin{omfigure}
\begin{tikzpicture}[econ]
\begin{axis}[omrecord, xmin=0.5, xmax=12.5, ymin=0, ymax=1, width=0.6\linewidth, height=0.4\linewidth,
    xlabel={Number of bidders $n$}, ylabel={Correction, share of $\bar\varepsilon$}, xtick={1,2,4,6,8,10,12}]
  \addplot[omDef, thick, mark=*, mark size=1.6pt, samples at={1,...,12}] {(x - 1)/(x + 1)};
\end{axis}
\end{tikzpicture}

{\small The amount by which the expected value of a common-value object falls short of the winner's signal, as a share of
the largest possible error. With one bidder winning reveals nothing; with more, the winner is more likely to hold a signal
far above the true value.}
\end{omfigure}

A rational bidder bids as if her signal were the highest, which lowers her bid, more so the more rivals there are. In a
private-values auction more rivals mean bidding closer to one's value; in a \omterm{def:b2:bayesian-games:common-value}{common-value auction} they also mean that
winning is worse news.

\begin{record}[Oil leases and the winner's curse]
% ledger: b2:thaler1988-oil, b2:thaler1988-jar
Engineers at an oil company argued in 1971 that winners of offshore oil leases earned poor returns because they were the
bidders who had most overestimated the oil (Capen, Clapp and Campbell 1971). The bids they reported were very dispersed: in
the 1969 sale of Alaska North Slope tracts the winning bids summed to 900 million dollars and the second-highest to 370
million, and the winning bid was at least twice the second on 77 per cent of tracts. A later study of leases in the Gulf
of Mexico found that, of 18 firms or consortia that bid often, 12 would have earned more by scaling down all their bids,
with a median optimal scaling of 0.68 (Hendricks, Porter and Boudreau 1987). In a classroom experiment, jars of coins worth
8 dollars, unknown to the bidders, were auctioned 48 times: the bidders' mean estimate of the value was 5.13 dollars, but
the mean winning bid was 10.01, a loss of 2.01 to the winner (Bazerman and Samuelson 1983). All three studies are as
summarised by Thaler (1988).
\end{record}

\section{A first look at perfect Bayesian equilibrium}\label{sec:b2:bayesian-games:pbe}

When a player with private information moves first, the others can learn from what she does. Take a job applicant
who is Good or Poor, with probability one half each, and knows which. She chooses whether to obtain a certificate,
which costs a Good applicant 1 and a Poor one 3. An employer sees whether she has the certificate, not her \omterm{def:b2:bayesian-games:type}{type}, and
hires her at a wage of 2 or rejects her. Hiring a Good applicant is worth 1 to the employer, a Poor one $-1$, rejecting
0.

\begin{omfigure}
\begin{tikzpicture}[econ, font=\footnotesize]
  \node[circle, draw, fill=white, inner sep=1.6pt] (N) at (0,0) {};
  \node[anchor=east] at (-0.15,0) {Nature};
  \node[circle, fill=black, inner sep=1.4pt] (G) at (0,1.6) {};
  \node[circle, fill=black, inner sep=1.4pt] (P) at (0,-1.6) {};
  \draw[thick] (N) -- node[right] {Good, $\tfrac12$} (G);
  \draw[thick] (N) -- node[right] {Poor, $\tfrac12$} (P);
  \node[above] at (G.north) {applicant};
  \node[below] at (P.south) {applicant};
  \foreach \nm/\x/\y in {GC/-3.4/1.6, PC/-3.4/-1.6, GN/3.4/1.6, PN/3.4/-1.6}
    \node[circle, fill=black, inner sep=1.4pt] (\nm) at (\x,\y) {};
  \draw[thick] (G) -- node[pos=0.6, above] {Certificate} (GC);
  \draw[thick] (P) -- node[pos=0.6, below] {Certificate} (PC);
  \draw[thick] (G) -- node[pos=0.6, above] {None} (GN);
  \draw[thick] (P) -- node[pos=0.6, below] {None} (PN);
  \draw[densely dashed] (GC) -- (PC);
  \draw[densely dashed] (GN) -- (PN);
  \node[anchor=west] at (-3.3,0) {employer};
  \node[anchor=east] at (3.3,0) {employer};
  \node[anchor=north west] at (-3.3,1.5) {$[\mu]$};
  \node[anchor=south west] at (-3.3,-1.5) {$[1 - \mu]$};
  \node[anchor=north east] at (3.3,1.5) {$[\mu']$};
  \node[anchor=south east] at (3.3,-1.5) {$[1 - \mu']$};
  \foreach \nm/\dx/\hy/\ry/\hp/\rp/\side in {
      GC/-1.5/2.4/0.8/{$(1, 1)$}/{$(-1, 0)$}/east, PC/-1.5/-0.8/-2.4/{$(-1, -1)$}/{$(-3, 0)$}/east,
      GN/1.5/2.4/0.8/{$(2, 1)$}/{$(0, 0)$}/west, PN/1.5/-0.8/-2.4/{$(2, -1)$}/{$(0, 0)$}/west} {
    \path (\nm) ++(\dx,0) coordinate (end);
    \draw[thick] (\nm) -- node[above, sloped] {hire} (end |- 0,\hy) node[anchor=\side] {\hp};
    \draw[thick] (\nm) -- node[below, sloped] {reject} (end |- 0,\ry) node[anchor=\side] {\rp};
  }
\end{tikzpicture}

{\small The certificate game. Nature draws the applicant's \omterm{def:b2:bayesian-games:type}{type}; she chooses a certificate or none; the employer, who sees
only that choice, hires or rejects. The dashed lines join the nodes the employer cannot tell apart, with her beliefs $\mu$
and $\mu'$ that the applicant is Good. Payoffs are the applicant's first and the employer's second.}
\end{omfigure}

\begin{definition}[Perfect Bayesian equilibrium]\label{def:b2:bayesian-games:pbe}
A \emph{perfect Bayesian equilibrium}\index{perfect Bayesian equilibrium} of a game in which players observe some
actions but not others' \omterm{def:b2:bayesian-games:type}{types} is a profile of strategies together with beliefs at every point where a player moves
without knowing the \omterm{def:b2:bayesian-games:type}{types}, such that each player's action is optimal given her beliefs and the others' strategies,
and beliefs follow from the strategies by Bayes' rule wherever that rule applies.
\end{definition}

In one equilibrium the Good applicant obtains the certificate and the Poor one does not. Bayes' rule then gives the
employer the beliefs $\mu = 1$ and $\mu' = 0$, so she hires with a certificate (1 is better than 0) and rejects without one. Each \omterm{def:b2:bayesian-games:type}{type} is
content: the Good applicant earns $2 - 1 = 1$ with a certificate and 0 without, the Poor applicant $2 - 3 = -1$ and 0. The
certificate separates the \omterm{def:b2:bayesian-games:type}{types}, though it teaches nothing.

There is another. Neither \omterm{def:b2:bayesian-games:type}{type} obtains the certificate, and the employer, believing $\mu' = \tfrac12$, expects $\tfrac12 - \tfrac12 = 0$ from
hiring and rejects. A certificate is never seen, so Bayes' rule says nothing about $\mu$; if the employer would believe a
certified applicant to be Poor with probability above one half, she would reject her too, and no applicant has a
reason to certify. Beliefs about actions that are never taken can sustain an equilibrium, and the signalling
models of \cref{ch:b2:signalling-and-screening} return to the question of which beliefs are reasonable.

\section{Exercises}\label{sec:b2:bayesian-games:exercises}

\begin{exercise}[$\star$]\label{exo:b2:bayesian-games:1}
Three bidders with values 30, 50 and 80 take part in a second-price auction. Who wins, and what does she pay? If their
values were drawn uniformly from $[0, 100]$, what would each bid in the first-price equilibrium, and who would win at what
price?
\end{exercise}

\begin{exercise}[$\star$]\label{exo:b2:bayesian-games:2}
Read the opening chart. How many of the eighteen subjects bid above the risk-neutral prediction with three bidders? With
six? What was the average bid per dollar of value in each case?
\end{exercise}

\begin{exercise}[$\star$]\label{exo:b2:bayesian-games:3}
Predict, then compute: the expected revenue of a first-price auction with values uniform on $[0, 1]$ and 2 bidders, then 10.
\end{exercise}

\begin{exercise}[$\star$]\label{exo:b2:bayesian-games:4}
What laboratory result would refute the prediction that bidders bid their values in a second-price auction? What
result would leave the prediction standing although bids differ from values?
\end{exercise}

\begin{exercise}[$\star\star$]\label{exo:b2:bayesian-games:5}
Redo the Cournot example of \cref{sec:b2:bayesian-games:bne} with probability 0.25 that firm 2's cost is high.
\end{exercise}

\begin{exercise}[$\star\star$]\label{exo:b2:bayesian-games:6}
Show by an example that bidding one's value is not a dominant strategy in a first-price auction.
\end{exercise}

\begin{exercise}[$\star\star$]\label{exo:b2:bayesian-games:7}
Thirty per cent of fishing grounds are rich. A survey reports ``rich'' for 80 per cent of rich grounds and for 20 per
cent of poor ones. What is the probability that a ground reported rich is rich? And one reported poor?
\end{exercise}

\begin{exercise}[$\star\star$]\label{exo:b2:bayesian-games:8}
In the jar experiment of the record box, what was the winner's average loss? Bidders' estimates of the value were
on average too low; why did the winners still pay too much?
\end{exercise}

\begin{exercise}[$\star\star$]\label{exo:b2:bayesian-games:9}
With \omterm{def:b2:expected-utility:crra}{constant relative risk aversion} $r$ the first-price bid is $(n - 1)v/(n - r)$. Find the bid per unit of value with three
bidders and $r = 0.5$. What $r$ would explain a slope of 0.75 with three bidders?
\end{exercise}

\begin{exercise}[$\star\star\star$]\label{exo:b2:bayesian-games:10}
Derive the first-price equilibrium for values uniform on $[0, \bar v]$ and $n$ bidders, assuming a linear bid function, and show
that no \omterm{def:b2:bayesian-games:type}{type} gains from any deviation.
\end{exercise}

\begin{exercise}[$\star\star\star$]\label{exo:b2:bayesian-games:11}
Show that the highest of $n$ independent values uniform on $[0, 1]$ has expectation $n/(n + 1)$ and the second-highest $(n -
1)/(n + 1)$, and deduce that both auction formats give the expected revenue $(n - 1)/(n + 1)$.
\end{exercise}

\begin{exercise}[$\star\star\star$]\label{exo:b2:bayesian-games:12}
Show that the highest of $n$ independent errors uniform on $[-\bar\varepsilon, \bar\varepsilon]$ has expectation $\bar\varepsilon(n - 1)/(n + 1)$.
\end{exercise}

\section{Problem: Brenland's fishing-licence auction}\label{sec:b2:bayesian-games:problem}

\begin{problem}[{Weekend problem --- the fisheries office of Brenland, an invented economy, auctions a one-year
licence to four fishing firms: how should they bid, which format raises more, and how should a firm bid when the
licence is worth the same to all}]
\label{pb:b2:bayesian-games:1}
Each firm's value for the licence is private and drawn uniformly between 0 and 800 thousand dalens. Firms are risk
neutral unless stated otherwise. Amounts are in thousands of dalens.

\medskip
\noindent\textbf{Part I --- Values and types.}
\begin{enumerate}
  \item What is a firm's \omterm{def:b2:bayesian-games:type}{type} in this auction?
  \item What is the probability that a given firm has the highest value?
  \item What is the expected highest value of the four?
  \item What is the expected second-highest value?
  \item What is the probability that a firm with value 600 has the highest value?
  \item Why is this a private-values auction?
\end{enumerate}

\medskip
\noindent\textbf{Part II --- First price.}
\begin{enumerate}[resume]
  \item Find the equilibrium bid function.
  \item What does a firm with value 600 bid?
  \item What is its expected payoff?
  \item What is the office's expected revenue?
  \item What would a firm with value 600 bid if all firms had \omterm{def:b2:expected-utility:risk-aversion}{risk aversion} $r = 0.5$?
  \item Why does no firm bid its value?
\end{enumerate}

\medskip
\noindent\textbf{Part III --- Second price.}
\begin{enumerate}[resume]
  \item What should each firm bid?
  \item What is the office's expected revenue?
  \item What is the expected payoff of a firm with value 600?
  \item The standard deviation of revenue is 98 in the first-price auction. In the second-price auction it is 160.
    Which format would an office that dislikes risk prefer?
  \item Which format raises more if the firms are averse to risk?
  \item Which format is easier for a firm to bid in, and why might that matter?
\end{enumerate}

\medskip
\noindent\textbf{Part IV --- A common value.} Now the licence gives access to a fishing ground whose catch is worth the
same amount $V$ to each firm. Each firm's survey gives a signal $s_i = V + \varepsilon_i$, with errors independent and uniform between
$-100$ and $100$.
\begin{enumerate}[resume]
  \item A firm's signal is 500. What is its expected value of the licence?
  \item What is its expected value if its signal turns out to be the highest of the four?
  \item A firm bids in a second-price auction as if the value were its signal. What does it expect to lose, if it wins
    with the highest signal, relative to bidding with the correction?
  \item By how much should a firm reduce its estimate with eight bidders?
  \item Why does the correction grow with the number of bidders?
  \item State the result: the expected revenue of both formats with private values, and the correction to a naive bid with
    a common value and four bidders.
\end{enumerate}
\end{problem}
